\documentclass[11pt,a4paper]{article}
\usepackage[utf8]{inputenc}
\usepackage{amsmath,amssymb,amsthm}
\usepackage[margin=1in]{geometry}

\newtheorem{theorem}{Theorem}[section]
\newtheorem{proposition}[theorem]{Proposition}
\newtheorem{remark}[theorem]{Remark}

\title{Negative Sectional Projections in $\mathfrak{su}(3)$ and $L=2^4, L=3^4$ Lattice Probes with Golden-Partitioned Conditional Stabilization and C99 Wilson Cross-Benchmarks}
\author{Mariano Panzano Caball\'e \\ Oasis Sovereign Architecture Research, Barcelona}
\date{November 2026}

\begin{document}
\maketitle

\begin{abstract}
We report finite-dimensional numerical probes on the Lie algebra $\mathfrak{su}(3)$ ($N=2000$ orthonormalized pairs) and on elementary $L=2^4$ and $L=3^4$ $\mathrm{SU}(3)$ lattices under adjoint link transport ($N_{\mathrm{cfg}}=200$ per lattice size). Across all three regimes, subtracting the orthogonal commutator contribution modeling quotient/horizontal curvature loss reverses the sign of the mean sectional quantity: $\bar{K}_{\mathrm{proj}} = -0.440979$ ($\min = -0.958835$) on pointwise $\mathfrak{su}(3)$, $\bar{K}_{\mathrm{proj}} = -0.396026$ ($\min = -0.875473$) on $L=2^4$, and $\bar{K}_{\mathrm{proj}} = -0.353224$ ($\min = -0.775408$) on $L=3^4$. We pair these measurements with a deterministic golden-ratio partition ($\varphi^{-1} \approx 0.618034$, $\varphi^{-2} \approx 0.381966$), cross-benchmark them against C99 Cabibbo--Marinari Wilson Monte Carlo observables on $L=8^4$ and $L=12^4$ lattices ($\beta \in [5.70, 6.10]$), and provide a formal Lean~4.36.0 conditional stabilization proof (\texttt{0 sorry}). These finite probes and conditional bounds do not establish unconditional continuum existence or a mass gap on $\mathbb{R}^4$ or $\mathcal{S}'(\mathbb{R}^4)$.
\end{abstract}

\section{Numerical Observations on $\mathfrak{su}(3)$, $L=2^4$, and $L=3^4$}
For Hilbert--Schmidt orthonormal generators $X, Y \in \mathfrak{su}(3)$, the intrinsic bi-invariant sectional quantity satisfies
\[
K_{\mathrm{int}}(X,Y) = \frac{1}{4}\|[X,Y]\|_{\mathrm{HS}}^2 \ge 0,
\]
while subtracting the orthogonal commutator contribution $K_{\mathrm{sub}}(X,Y) = \frac{3}{4}\|([X,Y])_\perp\|_{\mathrm{HS}}^2$ yields $K_{\mathrm{proj}} = K_{\mathrm{int}} - K_{\mathrm{sub}}$. Table~\ref{tab:probes} summarizes all three regimes.

\begin{table}[h]
\centering
\begin{tabular}{lcccc}
\hline
Probe Regime & Samples & $\bar{K}_{\mathrm{int}}$ & $\bar{K}_{\mathrm{proj}}$ & $\min K_{\mathrm{proj}}$ \\
\hline
Pointwise $\mathfrak{su}(3)$ algebra & $2000$ & $0.220498$ & $-0.440979$ & $-0.958835$ \\
Adjoint-transported $L=2^4$ lattice & $200$ & $0.198013$ & $-0.396026$ & $-0.875473$ \\
Adjoint-transported $L=3^4$ lattice & $200$ & $0.176612$ & $-0.353224$ & $-0.775408$ \\
\hline
\end{tabular}
\caption{Comparison of intrinsic and orthogonally projected sectional quantities.}
\label{tab:probes}
\end{table}

\section{Golden Partition and Conditional Margin Bound}
Using $\varphi = \frac{1+\sqrt{5}}{2}$, the complementary weights $\varphi^{-1} \approx 0.618034$ and $\varphi^{-2} \approx 0.381966$ define a deterministic parameter split for bounded stabilization checks.
\begin{proposition}[Conditional Bound Propagation]
Suppose a regulated model satisfies an explicit effective curvature-entropy margin $\kappa_{\mathrm{eff}} > 0$ and a positive effective string tension floor $\chi > 0$ under stated stabilization assumptions. Under the formal chain verified in Lean~4.36.0 (\texttt{0 sorry}), a strictly positive combined floor follows conditionally from those premises.
\end{proposition}

\section{Cross-Benchmark with C99 Wilson Monte Carlo ($L=8^4, 12^4$)}
To test finite-volume scaling under the Wilson plaquette action, C99 Cabibbo--Marinari runs were evaluated across four golden-partitioned couplings $\beta \in \{5.7000, 5.8673, 5.9706, 6.1000\}$ (Table~\ref{tab:wilson}).

\begin{table}[h]
\centering
\begin{tabular}{ccccc}
\hline
Lattice & $\beta$ & $\langle P \rangle$ & $W(2,2)$ & $\chi(2,2)$ \\
\hline
$8^4$ & $5.7000$ & $0.57071$ & $0.17223$ & $0.26832$ \\
$8^4$ & $5.9706$ & $0.59839$ & $0.20471$ & $0.23128$ \\
$8^4$ & $6.1000$ & $0.61517$ & $0.22503$ & $0.21746$ \\
$12^4$ & $5.7000$ & $0.57591$ & $0.18188$ & $0.24112$ \\
$12^4$ & $5.9706$ & $0.60700$ & $0.22074$ & $0.20649$ \\
$12^4$ & $6.1000$ & $0.61747$ & $0.23308$ & $0.19785$ \\
\hline
\end{tabular}
\caption{Mean plaquette $\langle P \rangle$, $2\times 2$ Wilson loop, and Creutz ratio $\chi(2,2)$.}
\label{tab:wilson}
\end{table}

\begin{remark}[Analytic Scope and Open Continuum Barrier]
At $40\text{--}60$ thermalization/measurement sweeps, direct $0^{++}$ correlator mass ratios exhibit statistical noise (\texttt{NaN}) and require higher statistics and smearing. Furthermore, pointwise $\mathfrak{su}(3)$ sampling, $L=2^4, 3^4$ link-transport probes, and finite-lattice Wilson runs do not establish reflection positivity in the continuum limit $a \to 0, L \to \infty$ or an unconditional mass gap on $\mathbb{R}^4$ or $\mathcal{S}'(\mathbb{R}^4)$.
\end{remark}

\end{document}
